\documentclass[10pt,a4paper]{amsart}
\usepackage[margin=19mm]{geometry}
\usepackage{amsmath,amssymb,mathtools}
\usepackage[scr=rsfs,cal=euler]{mathalfa}
\usepackage{xfrac}
\usepackage[unicode,hidelinks,bookmarks=false]{hyperref}
\newcommand{\ie}{i.e.\ }
\newcommand{\cf}{cf.\ }
\newcommand{\resp}{respectively\ }
\newcommand{\Subsection}{\subsection}
\providecommand{\nobreakdash}{\nobreak\hbox{-}}
\setlength{\emergencystretch}{2em}
\multlinegap0pt
\usepackage{bm}
\usepackage{bbm}
\usepackage{dsfont}
\usepackage{xspace}
\usepackage{cancel}
\usepackage{mathtools}
\usepackage{amsthm}
\makeatletter
\@ifundefined{theorem}{\newtheorem{theorem}{Theorem}}{}
\makeatother
\title[Coordinate-View Confusability Graphs and Matroid Rank Certif]{Coordinate-View Confusability Graphs and Matroid Rank Certificates}
\author{Tristan Simas}
\date{Source: arXiv:2602.23520v10 (CC BY 4.0)}
\begin{document}
\maketitle

\noindent\emph{Source and licence.} Coordinate-View Confusability Graphs and Matroid Rank Certificates. Version arXiv:2602.23520v10 (CC BY 4.0), distributed under the Creative Commons Attribution 4.0 International licence, as recorded in the arXiv licence metadata for this exact version and shown on the abstract page https://arxiv.org/abs/2602.23520v10. Full rights notice: licenses/arxiv-2602.23520-CC-BY-4.0.txt.\par\smallskip

\noindent\textit{Editorial scope.} Counted unit: the Theorem whose source label is \texttt{thm:intersection-closure} in the article (source0.tex lines 1359--1375, source label \texttt{thm:intersection-closure}), delivered together with the article's own complete original proof of it (source0.tex lines 1377--1398, one \texttt{proof} environment of 22 lines). No mathematics is written, repaired or re-derived: statement and proof are the article's own text line for line. Required context retained uncounted: none. Every definition, display and label of those lines is retained unchanged. Omitted from this package and not redistributed: 1--1358, 1376--1376, 1399--1573. Nothing inside a retained excerpt is omitted or altered: every retained body is delivered exactly as the article's lines are written, so no omission receipt of any kind accompanies this package. Citations: the retained excerpts carry no citation command at all, so no bibliography entry of the article is transcribed and no reference is redistributed. Reference adaptation: none was needed and none was made; every reference inside the delivered unit resolves inside it, so no unresolved reference or citation is produced. Printed slips kept literal: no printed word, symbol, hypothesis or number of the counted statement or of its proof was altered. Layout and class: the article's own class is \texttt{siamart} with packages including inputenc, fontenc, amsmath, amssymb, graphicx, xcolor, fancyvrb, url, listings, algorithm, algpseudocode, booktabs, longtable, array; the wrapper is the lane's pinned \texttt{amsart} editorial wrapper at 10pt on A4, which restates the theorem-like environments the retained text needs and adds the article's own macro definitions that the retained text needs (none), with extra wrapper packages (bm, bbm, dsfont, xspace, cancel, mathtools). The delivered numbering is the unit's own. Assets: no retained span contains a figure environment, a graphics-inclusion command, a PDF-page inclusion, a table or a TikZ picture; no image file is copied into this package, the package declares no asset dependency and no image file is redistributed with it. Licence: version arXiv:2602.23520v10 of this article is distributed under the Creative Commons Attribution 4.0 International licence, as recorded in the arXiv licence metadata for this exact version and shown on the abstract page https://arxiv.org/abs/2602.23520v10. A full rights notice is delivered at licenses/arxiv-2602.23520-CC-BY-4.0.txt. Characters: every retained excerpt is pure ASCII, so the delivered engine, XeLaTeX with its default Latin Modern text font, reports no missing character.
\begin{theorem}[Cluster/non-transitive dichotomy]\label{thm:intersection-closure}
For deterministic coordinate-subset views on the full tuple space over a nontrivial finite alphabet, meaning \( |A_i|\ge 2 \) for each coordinate,
\[
\begin{aligned}
G_{\mathrm{conf}}\text{ is transitive}
&\Longleftrightarrow
\mathcal U_{\mathcal V}\text{ is intersection-closed}\\
&\Longleftrightarrow
\mathcal V\text{ is meet-witnessed}.
\end{aligned}
\]
If these equivalent conditions hold, \(G_{\mathcal V}\) is a cluster graph and
\[
\Theta(G_{\mathcal V})=\vartheta_\infty(G_{\mathcal V})=\log |\Pi_{\mathrm{cc}}|.
\]
If they fail, there exist distinct states \(x,y,z\) with \(x\sim y\), \(y\sim z\), and \(x\not\sim z\).
\end{theorem}

\begin{proof}
Meet-witnessing and intersection closure are equivalent for the upward family. If \(R,R'\in\mathcal U_{\mathcal V}\), choose admissible \(S\subseteq R\) and \(S'\subseteq R'\). A meet witness \(W\subseteq S\cap S'\) gives \(R\cap R'\in\mathcal U_{\mathcal V}\). Conversely, applying intersection closure to two admissible views \(S,S'\) shows that \(S\cap S'\in\mathcal U_{\mathcal V}\), hence some admissible \(W\subseteq S\cap S'\) exists.

Intersection closure implies transitivity because \(x\sim y\) and \(y\sim z\) give \(\operatorname{Agr}(x,y),\operatorname{Agr}(y,z)\in\mathcal U_{\mathcal V}\); their intersection is contained in \(\operatorname{Agr}(x,z)\), and upward closure gives \(x\sim z\).

For the converse, choose \(R,R'\in\mathcal U_{\mathcal V}\) with \(R\cap R'\notin\mathcal U_{\mathcal V}\). Partition
\[
A=R\cap R',\quad B=R\setminus R',\quad C=R'\setminus R,\quad D=[d]\setminus(R\cup R').
\]
Using two symbols \(0,1\), define states by
\[
\begin{array}{c|cccc}
 & A & B & C & D\\ \hline
x & 0 & 0 & 1 & 0\\
y & 0 & 0 & 0 & 0\\
z & 0 & 1 & 0 & 1 .
\end{array}
\]
Then \(\operatorname{Agr}(x,y)\supseteq R\), \(\operatorname{Agr}(y,z)=R'\), and \(\operatorname{Agr}(x,z)=R\cap R'\notin\mathcal U_{\mathcal V}\). Thus \(x\sim y\), \(y\sim z\), and \(x\not\sim z\). The sets \(B\) and \(C\) are nonempty, since otherwise \(R\cap R'\) would equal one of \(R,R'\), so the three states are distinct.

In the transitive case, a finite graph with transitive adjacency is a disjoint union of cliques. One representative from each connected component is independent, and the component partition colors the complement. The standard sandwich \(\alpha(G)\le \vartheta(G)\le\chi(\overline G)\), together with multiplicativity of \(\vartheta\), gives the displayed capacity formula.
\end{proof}

\end{document}
